\documentclass[prx,amsmath,amssymb,superscriptaddress,longbibliography]{revtex4-2}
% \documentclass[10pt,prl,superscriptaddress,longbibliography]{revtex4-2}
\usepackage{xcolor}
\usepackage{bm}
\usepackage{amsmath}
\usepackage[colorlinks=true,linkcolor=blue,anchorcolor=blue,citecolor=blue,urlcolor=blue]{hyperref}


\usepackage{booktabs}
\usepackage{mathrsfs}
\makeatletter
\def\ltxu@dotsep{0.5em}
\makeatother
\newcommand{\Tr}{\operatorname{Tr}}
\newcommand{\Rea}{\operatorname{Re}}
\newcommand{\Ima}{\operatorname{Im}}
\newcommand{\dd}{\mathrm{d}}
\newcommand{\ii}{\mathrm{i}}
\newcommand{\symab}{\mathop{\mathrm{Sym}}\nolimits_{ab}}
\newcommand{\code}[1]{\texttt{\detokenize{#1}}}
\allowdisplaybreaks[2]
\begin{document}
\title{Personally implemented formulae in WannierBerri\\
Working handbook: geometry, covariant derivatives, and interpolation}
\author{Jiayu Li}
\email{jyli93@hku.hk}
\affiliation{Department of Physics, Southern University of Science and Technology, Shenzhen, China}
\affiliation{New Cornerstone Science Laboratory, Department of Physics, The University of Hong Kong, Hong Kong, China}
\date{Working copy, 5 October 2026}
\maketitle



This working copy consolidates the results in \code{formulae.tex},
\code{Wannier.tex}, and \code{QGD.tex}.
It establishes common conventions for the subsequent reconciliation of the
v1.7 and v1.8 calculators. Formula identities, representation assumptions,
and implementation status are stated separately.
The original source notes remain unchanged. Source corrections and unfinished
results are recorded in Appendix~\ref{app:corrections}.

% A compact contents list keeps the overview on the title page.
\begingroup
\small
\makeatletter
\def\l@f@section{\addvspace{0.4em}}
\makeatother
\tableofcontents
\endgroup
\newpage

\section{Basic ingredients}
\label{sec:basic}

\subsection{Notation, units, and scope}
\label{sec:notation}
We use Cartesian wave-vector derivatives
$\partial_a=\partial/\partial k_a$ and $\epsilon_{xyz}=+1$.
Repeated Cartesian indices are summed only in explicitly contracted expressions;
free indices remain free. The index convention throughout this handbook is
\begin{equation}
 \alpha,\beta,\gamma\in\{A,B\},\qquad
 m,n\in A,\qquad l,p\in B,\qquad
 j,k\in\mathrm{Full}=A\oplus B.
 \label{eq:index-spaces}
\end{equation}
Greek indices label blocks; $A,B$ themselves denote the chosen blocks.
Latin $m,n,l,p,j,k$ label states, whereas $a,b,c,d,e,\ldots$ label
Cartesian directions. Band and block indices are subscripts; Cartesian
components and derivative indices of matrix elements are superscripts,
as in $X_{mn}^{ab:d}$ and $\Delta_{jk}^{,d}$.
Differential operators retain their usual subscripts, e.g.\ $\partial_a$
and $\nabla_a$. Primes on band indices do not change their index space.
In the matrix-element formulae, every primed band index is implicitly summed
once over its assigned space:
$m',n',m'',n''\in A$, $l',p',l'',p''\in B$, and
$j',k',j'',k''\in\mathrm{Full}$.
Thus primed dummy indices do not require an explicit summation sign.
Free indices $m,n$ are not summed unless a sum or trace is written explicitly.
For example,
\begin{equation}
 (X_{AB}Y_{BA})_{mn}
 =\sum_{l\in B}X_{ml}Y_{ln}\equiv X_{ml'}Y_{l'n},\qquad
 (X_{AA}Y_{AA})_{mn}=X_{mn'}Y_{n'n}.
 \label{eq:component-sums}
\end{equation}
The first product sums over $B$; the second sums over $A$.
For full-space products and commutators, the same convention gives
\begin{equation}
 (XY)_{jk}=X_{jj'}Y_{j'k},\qquad
 [X,Y]_{jk}=X_{jj'}Y_{j'k}-Y_{jj'}X_{j'k}.
 \label{eq:full-product-components}
\end{equation}
Explicit matrix elements accompany the compact block identities below.
The blocks are separated by a nonzero gap at the wave vector under
consideration; degeneracies within a block are allowed.
An isolated-band formula selects $A=\{n\}$, with its complement in $B$;
unrestricted dummy sums use $j,k$, with exclusions written explicitly.
Here Full means the full band space of the Hamiltonian being considered.
For a finite Wannier model it is the represented space, not all physical
states; external contributions are treated separately below.
In Wannier-gauge matrices, $j,k$ label basis functions spanning that same
represented space, rather than energy eigenstates.
The spatial dimension is denoted by $d_{\mathrm{BZ}}$, independently of
Cartesian index $d$.
\begin{equation}
 \Delta_{jk}=\varepsilon_j-\varepsilon_k,\qquad
 \int[\dd\bm{k}]=\int_{\mathrm{BZ}}\frac{\dd^{d_{\mathrm{BZ}}} k}{(2\pi)^{d_{\mathrm{BZ}}}},\qquad
 e>0,\quad q_{\mathrm{electron}}=-e.
 \label{eq:conventions}
\end{equation}
The geometric expressions use energies and wave vectors, without silently
setting $\hbar=1$. In particular,
\begin{equation}
 v_n^a=\frac{1}{\hbar}\partial_a\varepsilon_n,\qquad
 E_n^{a}=\partial_a\varepsilon_n.
 \label{eq:velocity-units}
\end{equation}
If energy is in eV and length in \AA, $E_n^a$ is in eV\,\AA;
it is not yet a velocity in m/s. The conversion is
$v_n^a=(1\,\mathrm{eV\ in\ J})(1\,\text{\AA\ in m})E_n^a/\hbar$.

\begin{center}
\setlength{\tabcolsep}{6pt}
\begin{tabular}{lll}
\toprule
Symbol & Meaning & Relation to the source notes\\
\midrule
$Q_n^{ab}$ & quantum geometric tensor (QGT) & QG; called QMT in \code{formulae.tex}\\
$G_n^{ab}$ & single-band quantum metric & uppercase throughout this handbook\\
$\widetilde G_{mn}^{ab}$ & subspace quantum metric & matrix elements in the active space\\
$\mathcal{G}_n^{ab}$ & Berry connection polarizability (BCP) & original normalization retained\\
$D^a$ & $U^\dagger\partial_a U$ & coefficient-space connection\\
$\bar A^a$ & $U^\dagger A^{W,a}U$ & rotated Wannier connection\\
$M^a$ & $\bar A^a+\ii D^a$ & auxiliary connection in BCP formulae\\
$\widetilde{\mathcal E}^{ab}$ & energy-weighted geometric matrix & renamed from $\widetilde G$ in \code{Wannier.tex}\\
$\mathfrak m^a$ & interband magnetic kernel & separates magnetic $M$ from BCP $M$\\
\bottomrule
\end{tabular}
\end{center}

{
\subsection{Berry connection}

Single band Berry connection (Abelian Berry connection):
\begin{equation}
    A_n^a(\bm{k}) = i \langle u_n(\bm{k})| \partial_{k_a} u_n(\bm{k}) \rangle \equiv i \langle u_n| \partial_{a} u_n \rangle 
\label{eq:berry-connection}\end{equation}
Matrix Berry connection (with interband elements for $j\ne k$):
\begin{equation}
    A_{jk}^a(\bm{k}) = i \langle u_j(\bm{k})| \partial_{k_a} u_k(\bm{k}) \rangle \equiv i \langle u_j | \partial_{a} u_k \rangle 
\end{equation}

}
The matrix Berry connection is Hermitian, $(A^a)^\dagger=A^a$.
Its full matrix transforms inhomogeneously under a wave-vector-dependent
change of basis. Off-block elements transform covariantly under independent
rotations within $A$ and $B$.

\subsection{Quantum geometric tensor (QGT)}

Quantum geometry tensor: 
\begin{align}
     Q_{n}^{ab}(\bm{k}) &=  \langle{\partial_{k_a} u_n(\bm{k})}|(1 - |u_n(\bm{k})\rangle\langle{u_n(\bm{k})}|)|{\partial_{k_b} u_n(\bm{k})}\rangle \nonumber \\
     &\equiv \langle \partial_{a} u_n |(1 - |u_n\rangle\langle u_n |)|{\partial_{b} u_n}\rangle  \nonumber  \\
     &= \sum_{j(\neq n)}  \langle \partial_a u_n | u_j \rangle\langle u_j | \partial_b u_n \rangle \nonumber \\
     &= \sum_{j(\neq n)}  A_{nj}^a A_{jn}^b \nonumber \\
     &= G_n^{ab}(\bm{k}) - \frac{i}{2}\Omega_n^{ab}(\bm{k}) 
\label{eq:qgt}\end{align}
Quantum metric tensor: 
\begin{equation}
    G_{n}^{ab}(\bm{k}) = \Rea Q_{n}^{ab}(\bm{k}) = \Rea \sum_{j(\neq n)}  A_{nj}^a A_{jn}^b
\label{eq:metric-curvature}\end{equation}
Berry curvature tensor:
\begin{equation}
    \Omega_{n}^{ab}(\bm{k}) = -2\Ima Q_{n}^{ab}(\bm{k}) = -2 \Ima \sum_{j(\neq n)}  A_{nj}^a A_{jn}^b
\end{equation}
Berry curvature vector:
\begin{equation}
    \Omega_{n}^{a}(\bm{k}) = \frac{1}{2}\epsilon^{abc}\Omega_{n}^{bc}(\bm{k}) = -\epsilon^{abc} \Ima \sum_{j(\neq n)}  A_{nj}^b A_{jn}^c = \epsilon^{abc} \partial_b A_n^c
\label{eq:curvature-vector}\end{equation}

The band sum in Eq.~\eqref{eq:qgt} requires a complete complement to $n$.
A finite-model band sum gives the geometry of that model; missing states in a
Wannier representation require the external terms of Sec.~\ref{sec:interpolation}.
The metric is symmetric and positive semidefinite, whereas
$\Omega_n^{ab}=-\Omega_n^{ba}$.

{
\subsection{Berry connection polarizability (BCP)}

Berry connection polarizability~\cite{GaoPRL2014FieldInducedPositional}: 
\begin{equation}
    \mathcal{G}_n^{ab}(\bm{k}) = \frac{1}{2} \sum_{j(\neq n)}  \frac{A_{nj}^a A_{jn}^b + A_{nj}^b A_{jn}^a}{\varepsilon_n - \varepsilon_j} = \Rea \sum_{j(\neq n)}  \frac{A_{nj}^a A_{jn}^b}{\varepsilon_n - \varepsilon_j}
\label{eq:bcp-single}\end{equation}

}
The normalization in Eq.~\eqref{eq:bcp-single} is retained from the original notes.
Any response formula imported from another convention must first be converted
to this definition; a factor of two must not be inferred from the name BCP alone.
The units of $Q$, $G$, and $\Omega$ are length squared; those of
$\mathcal{G}$ are length squared divided by energy.

\section{Covariant derivatives and Wannier conventions}
\label{sec:covariant}

\subsection{Gauge dictionary and scope}
\label{sec:gauge-dictionary}
Here a gauge specifies a wave-vector-dependent orthonormal frame for a fixed
space of electronic states. It concerns Bloch-state and basis choices,
rather than the electromagnetic gauge of an applied field.
The relevant choices are summarized below; their transformation laws are
specified after the Wannier and Hamiltonian conventions.

\begin{center}
\setlength{\tabcolsep}{5pt}
\renewcommand{\arraystretch}{1.15}
\begin{tabular}{p{0.19\linewidth}p{0.35\linewidth}p{0.37\linewidth}}
\toprule
Choice & Defining frame or transformation & Meaning in this handbook\\
\midrule
Wannier gauge & $\{|u_j^W\rangle\}$ from a chosen set of Wannier functions
& Smooth interpolation basis; $\mathcal H^W$ is generally not diagonal.\\
Hamiltonian gauge & $|u_n\rangle=\sum_j|u_j^W\rangle U_{jn}$
& Energy eigenbasis; $U^\dagger\mathcal H^W U=\varepsilon$.\\
Block gauge & $W=\mathrm{diag}(W_A,W_B)$
& Independent rotations within fixed active and complementary subspaces;
the Hamiltonian remains block diagonal.\\
Single-band phase & $|u_n\rangle\mapsto e^{\ii\phi_n}|u_n\rangle$
& Residual $U(1)$ freedom of a nondegenerate energy eigenstate.\\
Parallel transport & A selected connection vanishes along a chosen path
& A local gauge-fixing condition; the connection and transport direction must be specified.\\
Wannier-center convention & Inclusion or omission of $e^{\ii\bm k\cdot\bm t_j}$
& A particular rephasing of the Wannier basis, requiring consistent changes to all matrix data.\\
\bottomrule
\end{tabular}
\end{center}
The Hamiltonian and Wannier gauges describe the same represented space.
Changing the disentanglement window or the represented subspace changes the
model or projector, and is not merely a gauge transformation.
No maximally localized gauge is required by the algebra below; localization
is a construction choice for the Wannier functions.

\subsection{Wannier and Hamiltonian gauges}
\label{sec:WH-gauges}
With Wannier centers $\bm t_j$, define
\begin{align}
 |u_{k\bm k}^{W}\rangle
 &=\sum_{\bm R}e^{\ii\bm k\cdot(\bm R+\bm t_k-\bm r)}|\bm Rk\rangle,
 \\
 \mathcal H_{jk}^{W}(\bm k)
 &=\sum_{\bm R}e^{\ii\bm k\cdot\bm R_{jk}}
 \langle\bm0 j|\hat{\mathcal H}|\bm Rk\rangle,
 \qquad \bm R_{jk}=\bm R+\bm t_k-\bm t_j,
 \label{eq:fourier}\\
 \mathcal H^W U&=U\varepsilon,\qquad
 |u_n\rangle=\sum_k|u_k^W\rangle U_{kn},\qquad
 \bar X=U^\dagger X^W U.
 \label{eq:gauges}
\end{align}
Normalization of the Bloch sums is understood.
Fourier derivatives follow by multiplying every real-space term by
$\ii R_{jk}^a$, with one factor for each derivative.

For the coefficient vectors (columns of $U$),
\begin{align}
 D^a&=U^\dagger\partial_a U,\qquad
 (D^a)^\dagger=-D^a,\qquad (D_{jk}^a)^*=-D_{kj}^a,
 \label{eq:D}\\
 D_{ln}^a&=\frac{\mathcal H_{ln}^{,a}}{\varepsilon_n-\varepsilon_l},
 \quad l\in B,\ n\in A,\qquad
 \mathcal H^{,a}=U^\dagger(\partial_a\mathcal H^W)U,
 \label{eq:D-gap}\\
 |\partial_a u_n\rangle
 &=\sum_k|\partial_a u_k^W\rangle U_{kn}+\sum_j|u_j\rangle D_{jn}^a,
 \\
 A^{H,a}&=\bar A^a+\ii D^a\equiv M^a,\qquad
 A_{jk}^{W,a}=\ii\langle u_j^W|\partial_a u_k^W\rangle.
 \label{eq:M}
\end{align}
The connection and rotation in matrix elements are
\begin{align}
 D_{jk}^a&=(U^\dagger)_{jj'}\partial_aU_{j'k},\nonumber\\
 \bar X_{jk}&=(U^\dagger)_{jj'}X_{j'k'}^WU_{k'k},\qquad
 M_{jk}^a=\bar A_{jk}^a+\ii D_{jk}^a.
 \label{eq:connection-components}
\end{align}

Diagonal entries of $D$ depend on the phase choice.
Equation~\eqref{eq:D-gap} is used only across separated blocks, never to divide
by a zero gap within a degenerate group.
In general, $D_{jk}^a\ne\langle u_j|\partial_a u_k\rangle$:
the latter also contains the derivative of the Wannier basis.

\subsection{Gauge transformations and invariant quantities}
\label{sec:gauge-transformations}
\paragraph{Rotations of the band frame.}
Hold the Wannier basis fixed and let $U'=UW$, where
$W=\mathrm{diag}(W_A,W_B)$ is unitary. Matrix multiplication below is over
band indices; all matrices and rotations can depend on $\bm k$.
For an operator matrix $X$ and the connections,
\begin{align}
 X'&=W^\dagger XW,\qquad
 D'{}^a=W^\dagger D^aW+W^\dagger\partial_aW,
 \label{eq:band-gauge-D}\\
 \bar A'{}^a&=W^\dagger\bar A^aW,\qquad
 A'{}^{H,a}=W^\dagger A^{H,a}W+\ii W^\dagger\partial_aW.
 \label{eq:band-gauge-A}
\end{align}
In full-space matrix elements these laws read
\begin{align}
 X'_{jk}&=(W^\dagger)_{jj'}X_{j'k'}W_{k'k},\nonumber\\
 D'{}_{jk}^a&=(W^\dagger)_{jj'}D_{j'k'}^aW_{k'k}
 +(W^\dagger)_{jj'}\partial_aW_{j'k},\nonumber\\
 \bar A'{}_{jk}^a&=(W^\dagger)_{jj'}\bar A_{j'k'}^aW_{k'k},\nonumber\\
 A'{}_{jk}^{H,a}&=(W^\dagger)_{jj'}A_{j'k'}^{H,a}W_{k'k}
 +\ii (W^\dagger)_{jj'}\partial_aW_{j'k}.
 \label{eq:band-gauge-components}
\end{align}

Thus the full $A^H$ and $D$ are connections, not covariant operator matrices.
Their off-block parts are covariant because the inhomogeneous term is block diagonal:
\begin{equation}
 X'_{AB}=W_A^\dagger X_{AB}W_B,\quad
 D'{}^a_{AB}=W_A^\dagger D^a_{AB}W_B,\quad
 M'{}^a_{AB}=W_A^\dagger M^a_{AB}W_B.
 \label{eq:offblock-gauge}
\end{equation}
With $W_{nl}=W_{ln}=0$, the same expressions reduce to
\begin{align}
 X'_{ml}&=(W^\dagger)_{mm'}X_{m'l'}W_{l'l},\nonumber\\
 D'{}_{ml}^a&=(W^\dagger)_{mm'}D_{m'l'}^aW_{l'l},\qquad
 M'{}_{ml}^a=(W^\dagger)_{mm'}M_{m'l'}^aW_{l'l},\nonumber\\
 X'_{mn}&=(W^\dagger)_{mm'}X_{m'n'}W_{n'n},\qquad
 \mathcal H'_{mn}=(W^\dagger)_{mm'}\varepsilon_{m'}W_{m'n}.
 \label{eq:offblock-gauge-components}
\end{align}

This is the gauge freedom used by the block-covariant construction of
Ref.~\cite{Liu2023CovariantderivativesBerry}.
An arbitrary $W_A$ may mix different energies within the active group.
Then $\mathcal H'_{AA}=W_A^\dagger\varepsilon_A W_A$ need not be diagonal,
although the physical subspace and its spectrum are unchanged.
The rotated frame is an admissible block gauge, but is generally no longer
the Hamiltonian gauge.

To stay in the Hamiltonian gauge, a continuous rotation must preserve the
energy eigenspaces. Nondegenerate eigenstates retain independent phases;
an exactly degenerate eigenspace permits a full unitary rotation.
A numerical grouping of nearly degenerate bands does not make their distinct
energies exactly equal. Explicit denominators $1/(\varepsilon_j-\varepsilon_k)$
refer to the energy eigenbasis. In a general block frame their covariant
meaning is the inverse Sylvester operator of Eq.~\eqref{eq:sylvester}.

\paragraph{Single-band phase freedom.}
For $W_{jk}=\delta_{jk}e^{\ii\phi_j}$, Eqs.~\eqref{eq:band-gauge-D} and
\eqref{eq:band-gauge-A} imply
\begin{align}
 A'{}^a_j&=A_j^a-\partial_a\phi_j,\qquad
 D'{}^a_{jj}=D^a_{jj}+\ii\partial_a\phi_j,
 \label{eq:phase-diagonal}\\
 A'{}^a_{jk}&=e^{\ii(\phi_k-\phi_j)}A^a_{jk},\qquad j\ne k.
 \label{eq:phase-offdiagonal}
\end{align}
Therefore $A_{nj}^aA_{jn}^b$, $Q_n^{ab}$, $G_n^{ab}$,
$\Omega_n^{ab}$, and $\mathcal{G}_n^{ab}$ are phase invariant for isolated bands.
Their individual band labels do not define invariant quantities under
arbitrary mixing of a degenerate multiplet; use the corresponding subspace objects there.

\paragraph{Changing the Wannier basis while holding Bloch eigenstates fixed.}
Write the Wannier frame as a row of kets, so that
$|u^{W}\rangle\mathsf S$ denotes a linear combination of its columns.
For a smooth unitary $\mathsf S(\bm k)$ within the same represented space,
\begin{align}
 |u^{W'}\rangle&=|u^W\rangle\mathsf S,\qquad U'=\mathsf S^\dagger U,\qquad
 |u^{H'}\rangle=|u^H\rangle,\nonumber\\
 \mathcal H^{W'}&=\mathsf S^\dagger\mathcal H^W\mathsf S,\qquad
 A^{W',a}=\mathsf S^\dagger A^{W,a}\mathsf S
                         +\ii\mathsf S^\dagger\partial_a\mathsf S.
 \label{eq:Wannier-frame-change}
\end{align}
Equivalently, for each basis state and matrix element,
\begin{align}
 |u_j^{W'}\rangle&=|u_{j'}^W\rangle\mathsf S_{j'j},\qquad
 U'_{jn}=(\mathsf S^\dagger)_{jj'}U_{j'n},\nonumber\\
 \mathcal H_{jk}^{W'}&=
 (\mathsf S^\dagger)_{jj'}\mathcal H_{j'k'}^W\mathsf S_{k'k},\nonumber\\
 A_{jk}^{W',a}&=
 (\mathsf S^\dagger)_{jj'}A_{j'k'}^{W,a}\mathsf S_{k'k}
 +\ii(\mathsf S^\dagger)_{jj'}\partial_a\mathsf S_{j'k}.
 \label{eq:Wannier-frame-components}
\end{align}

Define $\mathcal K^a=U^\dagger(\partial_a\mathsf S)\mathsf S^\dagger U$.
In the unchanged physical band frame,
\begin{equation}
 D'{}^a=D^a-\mathcal K^a,\qquad
 \bar A'{}^a=\bar A^a+\ii\mathcal K^a,\qquad
 \bar A'{}^a+\ii D'{}^a=A^{H,a}.
 \label{eq:connection-cancellation}
\end{equation}
The compensating terms can also be read element by element:
\begin{align}
 \mathcal K_{jk}^a&=
 (U^\dagger)_{jj'}(\partial_a\mathsf S_{j'k'})(\mathsf S^\dagger)_{k'j''}U_{j''k},\nonumber\\
 D'{}_{jk}^a&=D_{jk}^a-\mathcal K_{jk}^a,\qquad
 \bar A'{}_{jk}^a=\bar A_{jk}^a+\ii\mathcal K_{jk}^a,\qquad
 \bar A'{}_{jk}^a+\ii D'{}_{jk}^a=A_{jk}^{H,a}.
 \label{eq:connection-cancellation-components}
\end{align}

The split into internal and external terms can consequently depend on the
Wannier representation. A complete expression for the same physical states
remains unchanged when every basis-derivative term is transformed consistently.
Dropping external terms is a model approximation, not a general gauge fixing.

\paragraph{Wannier-center phase convention.}
Let $|u_j^{W,0}\rangle$ omit the center phase and define
$\mathsf S_t=\mathrm{diag}(e^{\ii\bm k\cdot\bm t_j})$.
The convention in Eq.~\eqref{eq:fourier} satisfies
\begin{align}
 |u_j^{W,t}\rangle&=e^{\ii\bm k\cdot\bm t_j}|u_j^{W,0}\rangle,
 &\mathcal H^{W,t}&=\mathsf S_t^\dagger\mathcal H^{W,0}\mathsf S_t,
 \nonumber\\
 A^{W,t,a}&=\mathsf S_t^\dagger A^{W,0,a}\mathsf S_t-\mathsf t^a,
 &\mathsf t^a_{jk}&=t_j^a\delta_{jk}.
 \label{eq:center-gauge}
\end{align}
In particular, the center convention gives
\begin{align}
 \mathcal H_{jk}^{W,t}
 &=e^{\ii\bm k\cdot(\bm t_k-\bm t_j)}\mathcal H_{jk}^{W,0},\nonumber\\
 A_{jk}^{W,t,a}
 &=e^{\ii\bm k\cdot(\bm t_k-\bm t_j)}A_{jk}^{W,0,a}-t_j^a\delta_{jk}.
 \label{eq:center-gauge-components}
\end{align}

The Hamiltonian Fourier phase changes from $e^{\ii\bm k\cdot\bm R}$ to
$e^{\ii\bm k\cdot(\bm R+\bm t_k-\bm t_j)}$.
This is a special case of Eq.~\eqref{eq:Wannier-frame-change}, including the
connection shift. Changing only the hopping phase while keeping all position
and derivative matrices unchanged is inconsistent.
The resulting periodicity across reciprocal-lattice boundaries can differ
between frames; physical quantities are compared using the corresponding boundary transformations.

\paragraph{Covariance, invariance, and projectors.}
A covariant matrix changes its components as $X'_{AA}=W_A^\dagger X_{AA}W_A$.
Its trace is invariant, whereas its individual diagonal elements generally are not.
The physical projector $P=\sum_{n\in A}|u_n\rangle\langle u_n|$ is unchanged
by rotations confined to $A$. The matrices $\widetilde F$, $\widetilde G$,
and $\widetilde\Omega$ are covariant, and their subspace traces are invariant.
With occupations or other energy-dependent weights, the invariant object is
$\Tr_A[f(\mathcal H_{AA})X_{AA}]$; both factors must be rotated.
In matrix elements,
\begin{equation}
 \Tr_A X=\sum_{n\in A}X_{nn},\qquad
 \Tr_A[f(\mathcal H_{AA})X_{AA}]
 =\sum_{m,n\in A}[f(\mathcal H_{AA})]_{mn}X_{nm}.
 \label{eq:weighted-trace-components}
\end{equation}
Rotations mixing $A$ with $B$ change the chosen active projector and are outside
this fixed-subspace covariance statement.

\subsection{Comma, ordinary, and covariant derivatives}
\label{sec:derivative-conventions}
For a specified Wannier-gauge matrix $X^W$, the comma derivative means
\begin{equation}
 X^{,d}=U^\dagger(\partial_d X^W)U,\qquad
 \partial_d\bar X=X^{,d}+\bar X D^d-D^d\bar X.
 \label{eq:ordinary-derivative}
\end{equation}
The matrix-element version uses primed intermediate states in Full:
\begin{align}
 X_{jk}^{,d}&=
 (U^\dagger)_{jj'}(\partial_dX_{j'k'}^W)U_{k'k},\nonumber\\
 \partial_d\bar X_{jk}&=X_{jk}^{,d}
 +(\bar X_{jj'}D_{j'k}^d-D_{jj'}^d\bar X_{j'k}).
 \label{eq:ordinary-derivative-components}
\end{align}

Thus $X^{,d}$ is not the ordinary derivative of $\bar X$.
In a fixed basis for an abstract operator it becomes
$\langle j|\partial_d\hat X|k\rangle$.
Such an operator-derivative matrix is covariant; it is the ordinary derivative
of a transformed matrix that acquires derivatives of the gauge rotation.

Under $U\mapsto U\,\mathrm{diag}(W_A,W_B)$,
$X_{AB}\mapsto W_A^\dagger X_{AB}W_B$.
Using the coefficient-space connection, define~\cite{Liu2023CovariantderivativesBerry}
\begin{equation}
 X_{\alpha\beta}^{:d}=\partial_d X_{\alpha\beta}+D_{\alpha\alpha}^d X_{\alpha\beta}-X_{\alpha\beta}D_{\beta\beta}^d,
 \quad \alpha,\beta\in\{A,B\}.
 \label{eq:cov-derivative}
\end{equation}
For each choice of row and column space this is
\begin{align}
 X_{mn}^{:d}&=\partial_dX_{mn}+D_{mn'}^dX_{n'n}-X_{mn'}D_{n'n}^d,\nonumber\\
 X_{ml}^{:d}&=\partial_dX_{ml}+D_{mn'}^dX_{n'l}-X_{ml'}D_{l'l}^d,\nonumber\\
 X_{ln}^{:d}&=\partial_dX_{ln}+D_{lp'}^dX_{p'n}-X_{lm'}D_{m'n}^d,\nonumber\\
 X_{lp}^{:d}&=\partial_dX_{lp}+D_{ll'}^dX_{l'p}-X_{ll'}D_{l'p}^d.
 \label{eq:cov-derivative-components}
\end{align}

Here $\alpha,\beta$ label blocks, while the superscript $d$ labels a Cartesian direction.
The symbol $X$ in this equation denotes a matrix in the chosen band frame,
which may be the Hamiltonian gauge or a block rotation of it.
For a rotated Wannier matrix the resulting block rules are
\begin{align}
 X_{AB}^{:d}&=X_{AB}^{,d}-D_{AB}^dX_{BB}+X_{AA}D_{AB}^d,\label{eq:block-ab}\\
 X_{BA}^{:d}&=X_{BA}^{,d}-D_{BA}^dX_{AA}+X_{BB}D_{BA}^d,\\
 X_{AA}^{:d}&=X_{AA}^{,d}-D_{AB}^dX_{BA}+X_{AB}D_{BA}^d.
 \label{eq:block-aa}
\end{align}
In matrix elements, these interpolation rules become
\begin{align}
 X_{ml}^{:d}&=X_{ml}^{,d}-D_{ml'}^dX_{l'l}+X_{mn'}D_{n'l}^d,\nonumber\\
 X_{ln}^{:d}&=X_{ln}^{,d}-D_{lm'}^dX_{m'n}+X_{lp'}D_{p'n}^d,\nonumber\\
 X_{mn}^{:d}&=X_{mn}^{,d}-D_{ml'}^dX_{l'n}+X_{ml'}D_{l'n}^d,\nonumber\\
 X_{lp}^{:d}&=X_{lp}^{,d}-D_{lm'}^dX_{m'p}+X_{lm'}D_{m'p}^d.
 \label{eq:interpolation-derivative-components}
\end{align}
In particular, $\mathcal H_{mn}^{:d}=\mathcal H_{mn}^{,d}$,
$\mathcal H_{lp}^{:d}=\mathcal H_{lp}^{,d}$, and
$\mathcal H_{ml}^{:d}=\mathcal H_{ln}^{:d}=0$.

For the Hamiltonian, $\mathcal H_{AA}^{:d}=\mathcal H_{AA}^{,d}$ and
$\mathcal H_{AB}^{:d}=0$.
The colon derivative is covariant under band-frame rotations at fixed
Wannier representation:
\begin{equation}
 (X'_{\alpha\beta})^{:a}=W_\alpha^\dagger X_{\alpha\beta}^{:a}W_\beta.
 \label{eq:colon-transformation}
\end{equation}
For example,
\begin{equation}
 (X'_{mn})^{:a}=(W^\dagger)_{mm'}X_{m'n'}^{:a}W_{n'n},\qquad
 (X'_{ml})^{:a}=(W^\dagger)_{mm'}X_{m'l'}^{:a}W_{l'l}.
 \label{eq:colon-transformation-components}
\end{equation}

Its definition uses $D$, rather than the full physical Bloch connection.
To distinguish the latter construction, write
\begin{align}
 \nabla_a X_{\alpha\beta}
 &\equiv\partial_aX_{\alpha\beta}-\ii A^{H,a}_{\alpha\alpha}X_{\alpha\beta}
                           +\ii X_{\alpha\beta}A^{H,a}_{\beta\beta}
 \nonumber\\
 &=X_{\alpha\beta}^{:a}-\ii\bar A^a_{\alpha\alpha}X_{\alpha\beta}+\ii X_{\alpha\beta}\bar A^a_{\beta\beta}.
 \label{eq:physical-covariant-derivative}
\end{align}
Its matrix elements in the four index spaces are
\begin{align}
 (\nabla_aX)_{mn}&=X_{mn}^{:a}-\ii\bar A_{mn'}^aX_{n'n}+\ii X_{mn'}\bar A_{n'n}^a,\nonumber\\
 (\nabla_aX)_{ml}&=X_{ml}^{:a}-\ii\bar A_{mn'}^aX_{n'l}+\ii X_{ml'}\bar A_{l'l}^a,\nonumber\\
 (\nabla_aX)_{ln}&=X_{ln}^{:a}-\ii\bar A_{lp'}^aX_{p'n}+\ii X_{lm'}\bar A_{m'n}^a,\nonumber\\
 (\nabla_aX)_{lp}&=X_{lp}^{:a}-\ii\bar A_{ll'}^aX_{l'p}+\ii X_{ll'}\bar A_{l'p}^a.
 \label{eq:physical-derivative-components}
\end{align}

Both derivatives transform covariantly under $W_\alpha,W_\beta$, but they need not
have the same matrix elements. The physical connection $A^H$ is unchanged
by the compensated Wannier-basis change in Eq.~\eqref{eq:connection-cancellation};
$D$ alone is not. Their difference has zero trace for an $AA$ block,
so both give $\partial_a\Tr_A X$ after taking that trace.
The block interpolation identities in this handbook use the colon convention;
replacing $D$ by $-\ii A^H$ in those identities would require rederiving the comma terms.

\paragraph{Abelian derivative used in the isolated-band formulae.}
For a matrix element transforming as $X'_{jk}=e^{\ii(\phi_k-\phi_j)}X_{jk}$,
our coefficient-space phase derivative is
\begin{equation}
 X_{jk}^{;a}\equiv\partial_a X_{jk}
 +(D_{jj}^a-D_{kk}^a)X_{jk}.
 \label{eq:abelian-coefficient-derivative}
\end{equation}
The physical Bloch-state version is
\begin{equation}
 (\nabla_a X)_{jk}=\partial_aX_{jk}-\ii(A_j^a-A_k^a)X_{jk}
 =X_{jk}^{;a}-\ii(\bar A_{jj}^a-\bar A_{kk}^a)X_{jk}.
 \label{eq:abelian-physical-derivative}
\end{equation}
They coincide in the fixed-basis model of Sec.~\ref{sec:qgd}, where
$A^W=0$ and $A^H=\ii D$. In a general Wannier basis, they must be distinguished.
The semicolon in the BCP quotient formulae refers to
Eq.~\eqref{eq:abelian-coefficient-derivative}.

\paragraph{Parallel transport is a local condition.}
Along a path $\bm k(s)$, coefficient parallel transport of a nondegenerate
band means $D_{nn}^{\parallel}=\sum_a\dot k_aD_{nn}^a=0$.
A phase can be chosen to impose this condition along an open smooth path,
or to set the diagonal connection to zero at a specified point.
Physical Bloch parallel transport instead means
$A_n^{H,\parallel}=\ii\langle u_n|\dd u_n/\dd s\rangle=0$.
These choices differ when $\bar A_{nn}^{\parallel}\ne0$.
The label $\parallel$ denotes contraction with the path tangent, not a band index.
Neither choice removes interband $D_{jk}$ or $A_{jk}^H$ ($j\ne k$).
Nonzero curvature prevents a connection from being set to zero throughout a
whole patch; nontrivial holonomy can also obstruct a periodic gauge with
zero connection around a closed loop.
A block parallel-transport frame can similarly set its tangential intrablock
connection to zero, but may mix distinct energies and leave the Hamiltonian gauge.

\paragraph{Reading the subsequent formulae.}
Use the fixed Wannier convention for the colon derivatives and Fourier inputs.
Evaluate explicit energy denominators in the Hamiltonian gauge, using block
identities when rotating the active space.
Use isolated-band phase derivatives only where the relevant bands are nondegenerate.
The notation $\bar X$ means a matrix rotation, and $\widetilde X$ denotes a
projected construction; neither symbol defines an additional gauge.

\subsection{Trace, product, and chain rules}
The trace and product rules eliminate gauge-dependent intrablock connections:
\begin{align}
 \partial_d\Tr_A X&=\Tr_A X^{:d},
 \label{eq:trace}\\
 (X_{\alpha\beta}Y_{\beta\gamma})^{:d}&=X_{\alpha\beta}^{:d}Y_{\beta\gamma}+X_{\alpha\beta}Y_{\beta\gamma}^{:d}.
 \label{eq:product}
\end{align}
In components, representative $A$- and $B$-space contractions are
\begin{align}
 \partial_d\sum_{n\in A}X_{nn}&=\sum_{n\in A}X_{nn}^{:d},\nonumber\\
 (X_{mn'}Y_{n'n})^{:d}&=X_{mn'}^{:d}Y_{n'n}+X_{mn'}Y_{n'n}^{:d},\nonumber\\
 (X_{ml'}Y_{l'n})^{:d}&=X_{ml'}^{:d}Y_{l'n}+X_{ml'}Y_{l'n}^{:d}.
 \label{eq:product-components}
\end{align}

Equation~\eqref{eq:product} uses fixed blocks
$\alpha,\beta,\gamma\in\{A,B\}$; the product contracts state indices
inside the intermediate block $\beta$, with no implicit sum over blocks.
For $\widetilde f=f(\mathcal H)$, evaluation in the energy eigenbasis gives
\begin{equation}
 \widetilde f_{mn}^{:d}=\mathcal H_{mn}^{:d}
 \begin{cases}
 f'(\varepsilon_n),&\varepsilon_m=\varepsilon_n,\\[2pt]
 \displaystyle\frac{f(\varepsilon_m)-f(\varepsilon_n)}{\varepsilon_m-\varepsilon_n},
 &\varepsilon_m\ne\varepsilon_n.
 \end{cases}
 \label{eq:chain}
\end{equation}
The divided difference is essential when a block contains distinct energies.
At zero temperature, apply the smooth-function rule before taking the
appropriate distributional limit of the occupation.

For compact second derivatives, define the bilinear product expansion
\begin{align}
 \mathscr P^{d}(X,Y)&=X^{:d}Y+XY^{:d},\nonumber\\
 \mathscr P^{de}(X,Y)&=X^{:de}Y+X^{:d}Y^{:e}
                 +X^{:e}Y^{:d}+XY^{:de},
 \label{eq:second-product}\\
 X^{:de}&\equiv(X^{:d})^{:e}.
 \nonumber
\end{align}
For an intermediate state in $B$, its fully indexed second-order form is
\begin{equation}
 [\mathscr P^{de}(X_{AB},Y_{BA})]_{mn}
 =X_{ml'}^{:de}Y_{l'n}+X_{ml'}^{:d}Y_{l'n}^{:e}
  +X_{ml'}^{:e}Y_{l'n}^{:d}+X_{ml'}Y_{l'n}^{:de}.
 \label{eq:second-product-components}
\end{equation}
For an intermediate state in $A$, replace $l'\in B$ by $n'\in A$
in all four terms.

Matrix order is preserved. Ordered covariant derivatives need not commute;
ordinary derivatives of a smooth traced scalar do commute.

\section{Gauge-covariant geometry and Wannier interpolation}
\label{sec:interpolation}

\subsection{Projectors and geometric matrices}
Let $P=\sum_{n\in A}|u_n\rangle\langle u_n|$ and $Q=1-P$.
For physical Bloch states, $Q$ includes states outside a finite Wannier space.
The projected geometric matrices are
\begin{align}
 \widetilde F_{mn}^{ab}
 &=\langle\partial_a u_m|Q|\partial_b u_n\rangle,
 \label{eq:Fprojected}\\
 \widetilde{\mathcal E}^{ab}&=\tfrac12(\mathcal H_{AA}\widetilde F^{ab}
                            +\widetilde F^{ab}\mathcal H_{AA}),
 \label{eq:Gweighted}\\
 \widetilde H_{mn}^{ab}
 &=\langle\partial_a u_m|Q\hat{\mathcal H}Q|\partial_b u_n\rangle.
 \label{eq:Hprojected}
\end{align}
The energy-weighted matrix in Eq.~\eqref{eq:Gweighted} has elements
\begin{equation}
 \widetilde{\mathcal E}_{mn}^{ab}=\tfrac12\left(
 \mathcal H_{mn'}\widetilde F_{n'n}^{ab}
 +\widetilde F_{mn'}^{ab}\mathcal H_{n'n}\right).
 \label{eq:Gweighted-components}
\end{equation}

Each transforms as $\widetilde O\mapsto W_A^\dagger\widetilde O W_A$ and obeys
$\widetilde O_{mn}^{ab}=(\widetilde O_{nm}^{ba})^*$.
In the Hamiltonian gauge,
$\widetilde{\mathcal E}_{mn}^{ab}=(\varepsilon_m+\varepsilon_n)\widetilde F_{mn}^{ab}/2$.
The metric and curvature matrices are
\begin{equation}
 \widetilde G^{ab}=\tfrac12(\widetilde F^{ab}+\widetilde F^{ba}),\qquad
 \widetilde\Omega^c=\ii\epsilon_{cab}\widetilde F^{ab}.
 \label{eq:projected-geometry}
\end{equation}
Element by element,
\begin{equation}
 \widetilde G_{mn}^{ab}=\tfrac12(\widetilde F_{mn}^{ab}+\widetilde F_{mn}^{ba}),\qquad
 \widetilde\Omega_{mn}^c=\ii\epsilon_{cab}\widetilde F_{mn}^{ab}.
 \label{eq:projected-geometry-components}
\end{equation}
Here $\widetilde G$ denotes the metric; $\widetilde{\mathcal E}$ in
Eq.~\eqref{eq:Gweighted-components} remains the energy-weighted matrix.

For $A=\{n\}$ these recover Sec.~\ref{sec:basic}.
For a larger $A$, $\Tr_A\widetilde G$ is the metric of the whole subspace;
it is generally not $\sum_{n\in A}G_n$.

For a continuum Hamiltonian with kinetic energy $(\bm p+\hbar\bm k)^2/(2m_e)$
and a local, momentum-independent potential, the fourth matrix in
\code{Wannier.tex} is
\begin{equation}
 \widetilde S_{mn}^{ab}=\frac{\hbar^2}{m_e}\delta_{ab}\delta_{mn}
                  +2(\widetilde{\mathcal E}_{mn}^{ab}-\widetilde H_{mn}^{ab}).
 \label{eq:S-continuum}
\end{equation}
For $A=\{n\}$, $\Rea\widetilde S_{nn}^{ab}=\partial_a\partial_b\varepsilon_n$.
The free-electron term cannot be substituted for the bare Hessian of a
general tight-binding or nonlocal Hamiltonian; use Sec.~\ref{sec:energy} there.

\subsection{Internal and external terms}
Define Wannier-gauge matrices
\begin{align}
 F_{jk}^{W,ab}&=\langle\partial_a u_j^W|\partial_b u_k^W\rangle,\nonumber\\
 B_{jk}^{W,a}&=\ii\langle u_j^W|\hat{\mathcal H}|\partial_a u_k^W\rangle,\nonumber\\
 H_{jk}^{W,ab}&=\langle\partial_a u_j^W|\hat{\mathcal H}|\partial_b u_k^W\rangle.
 \label{eq:external-data}
\end{align}
The two-index $H^{W,ab}$ is an energy-weighted geometric matrix,
not $\partial_a\partial_b\mathcal H^W$.
Then
\begin{align}
 \widetilde F^{ab}_{\mathrm{int}}
 &=-D_{AB}^aD_{BA}^b,\label{eq:Fint}\\
 \widetilde F^{ab}_{\mathrm{ext}}
 &=\bar F_{AA}^{ab}-\bar A_{AA}^a\bar A_{AA}^b
    +\ii D_{AB}^a\bar A_{BA}^b+\ii\bar A_{AB}^a D_{BA}^b,\label{eq:Fext}\\
 \widetilde F^{ab}&=\widetilde F^{ab}_{\mathrm{int}}+
                    \widetilde F^{ab}_{\mathrm{ext}}.
\end{align}
The same internal and external terms, with every band index shown, are
\begin{align}
 \widetilde F_{mn,\mathrm{int}}^{ab}
 &=-D_{ml'}^aD_{l'n}^b,\nonumber\\
 \widetilde F_{mn,\mathrm{ext}}^{ab}
 &=\bar F_{mn}^{ab}-\bar A_{mn'}^a\bar A_{n'n}^b
 +\ii D_{ml'}^a\bar A_{l'n}^b+\ii\bar A_{ml'}^aD_{l'n}^b,\nonumber\\
 \widetilde F_{mn}^{ab}&=\widetilde F_{mn,\mathrm{int}}^{ab}
                         +\widetilde F_{mn,\mathrm{ext}}^{ab}.
 \label{eq:Finterpolation-components}
\end{align}
Thus the $DD$ and mixed products sum over $l'\in B$, while the
$\bar A\bar A$ subtraction sums over $n'\in A$.

Here $B$ is the complement of $A$ within the represented Wannier space;
$\bar F$ retains the additional basis-derivative information.
Keeping only $\widetilde F_{\mathrm{int}}$ corresponds to the fixed-basis
Hamiltonian geometry.

For $\widetilde H$, assuming the represented active states are eigenstates
of the physical Hamiltonian (or lie in the exact frozen subspace),
\begin{align}
 \widetilde H^{ab}_{\mathrm{int}}&=-D_{AB}^a\mathcal H_{BB}D_{BA}^b,\nonumber\\
 \widetilde H^{ab}_{\mathrm{ext}}
 &=\bar H_{AA}^{ab}-\bar A_{AA}^a\mathcal H_{AA}\bar A_{AA}^b
    +\ii D_{AB}^a\bar B_{BA}^b
    +\ii(\bar B^a)^\dagger_{AB}D_{BA}^b.
 \label{eq:Hinterpolation}
\end{align}
The corresponding matrix elements are
\begin{align}
 \widetilde H_{mn,\mathrm{int}}^{ab}
 &=-D_{ml'}^a\mathcal H_{l'p'}D_{p'n}^b,\nonumber\\
 \widetilde H_{mn,\mathrm{ext}}^{ab}
 &=\bar H_{mn}^{ab}-\bar A_{mn'}^a\mathcal H_{n'n''}\bar A_{n''n}^b
 +\ii D_{ml'}^a\bar B_{l'n}^b
 +\ii(\bar B_{l'm}^a)^*D_{l'n}^b.
 \label{eq:Hinterpolation-components}
\end{align}
In the Hamiltonian gauge the first line is
$\widetilde H_{mn,\mathrm{int}}^{ab}=-D_{ml'}^a\varepsilon_{l'}D_{l'n}^b$,
and the subtracted external product is
$-\bar A_{mn'}^a\varepsilon_{n'}\bar A_{n'n}^b$.

Outside that exact subspace, projection and disentanglement errors must be
assessed before interpreting the interpolation as a physical matrix element.

A useful sign check is
\begin{equation}
 \Omega_{A,\mathrm{int}}^{ab}
 =2\Ima\Tr_A(D_{AB}^aD_{BA}^b),\qquad
 \widetilde\Omega^z=\ii(\widetilde F^{xy}-\widetilde F^{yx}).
 \label{eq:curvature-check}
\end{equation}
Equivalently,
\begin{equation}
 \Omega_{A,\mathrm{int}}^{ab}
 =2\Ima\sum_{n\in A}\sum_{l\in B}D_{nl}^aD_{ln}^b,\qquad
 \widetilde\Omega_{mn}^z=\ii(\widetilde F_{mn}^{xy}-\widetilde F_{mn}^{yx}).
 \label{eq:curvature-check-components}
\end{equation}


\subsection{First and second geometric derivatives}
The internal first derivative is
\begin{equation}
 \widetilde F_{\mathrm{int}}^{ab:d}
 =-D_{AB}^{a:d}D_{BA}^b-D_{AB}^aD_{BA}^{b:d}.
 \label{eq:Ffirst}
\end{equation}
The external first derivative contains all seven product-rule terms:
\begin{align}
 \widetilde F_{\mathrm{ext}}^{ab:d}
 &=\bar F_{AA}^{ab:d}
 -\bar A_{AA}^{a:d}\bar A_{AA}^b-\bar A_{AA}^a\bar A_{AA}^{b:d}
 \nonumber\\
 &\quad+\ii D_{AB}^{a:d}\bar A_{BA}^b+\ii D_{AB}^a\bar A_{BA}^{b:d}
 +\ii\bar A_{AB}^{a:d}D_{BA}^b+\ii\bar A_{AB}^aD_{BA}^{b:d}.
 \label{eq:Fext-first}
\end{align}
For implementation, the explicit first-derivative matrix elements are
\begin{align}
 \widetilde F_{mn,\mathrm{int}}^{ab:d}
 &=-D_{ml'}^{a:d}D_{l'n}^b-D_{ml'}^aD_{l'n}^{b:d},\nonumber\\
 \widetilde F_{mn,\mathrm{ext}}^{ab:d}
 &=\bar F_{mn}^{ab:d}-\bar A_{mn'}^{a:d}\bar A_{n'n}^b
                        -\bar A_{mn'}^a\bar A_{n'n}^{b:d}\nonumber\\
 &\quad+\ii D_{ml'}^{a:d}\bar A_{l'n}^b
       +\ii D_{ml'}^a\bar A_{l'n}^{b:d}
       +\ii\bar A_{ml'}^{a:d}D_{l'n}^b
       +\ii\bar A_{ml'}^aD_{l'n}^{b:d}.
 \label{eq:Ffirst-components}
\end{align}
The derivative does not change the intermediate-state space: the two
$\bar A\bar A$ terms still contract $n'\in A$, not $l'\in B$.

All second-order terms are retained compactly by Eq.~\eqref{eq:second-product}:
\begin{align}
 \widetilde F_{\mathrm{int}}^{ab:de}
 &=-\mathscr P^{de}(D_{AB}^a,D_{BA}^b),\nonumber\\
 \widetilde F_{\mathrm{ext}}^{ab:de}
 &=\bar F_{AA}^{ab:de}-\mathscr P^{de}(\bar A_{AA}^a,\bar A_{AA}^b)
 \nonumber\\
 &\quad+\ii\mathscr P^{de}(D_{AB}^a,\bar A_{BA}^b)
       +\ii\mathscr P^{de}(\bar A_{AB}^a,D_{BA}^b).
 \label{eq:Fsecond}
\end{align}
Expanding the second derivative gives
\begin{align}
 \widetilde F_{mn,\mathrm{int}}^{ab:de}
 &=-D_{ml'}^{a:de}D_{l'n}^b-D_{ml'}^{a:d}D_{l'n}^{b:e}
   -D_{ml'}^{a:e}D_{l'n}^{b:d}-D_{ml'}^aD_{l'n}^{b:de},\nonumber\\
 \widetilde F_{mn,\mathrm{ext}}^{ab:de}
 &=\bar F_{mn}^{ab:de}
   -\bar A_{mn'}^{a:de}\bar A_{n'n}^b
   -\bar A_{mn'}^{a:d}\bar A_{n'n}^{b:e}\nonumber\\
 &\quad-\bar A_{mn'}^{a:e}\bar A_{n'n}^{b:d}
   -\bar A_{mn'}^a\bar A_{n'n}^{b:de}\nonumber\\
 &\quad+\ii D_{ml'}^{a:de}\bar A_{l'n}^b
   +\ii D_{ml'}^{a:d}\bar A_{l'n}^{b:e}
   +\ii D_{ml'}^{a:e}\bar A_{l'n}^{b:d}
   +\ii D_{ml'}^a\bar A_{l'n}^{b:de}\nonumber\\
 &\quad+\ii\bar A_{ml'}^{a:de}D_{l'n}^b
   +\ii\bar A_{ml'}^{a:d}D_{l'n}^{b:e}
   +\ii\bar A_{ml'}^{a:e}D_{l'n}^{b:d}
   +\ii\bar A_{ml'}^aD_{l'n}^{b:de}.
 \label{eq:Fsecond-components}
\end{align}

Consequently,
\begin{align}
 \widetilde G^{ab:d}&=\tfrac12(\widetilde F^{ab:d}+\widetilde F^{ba:d}),
 &\widetilde\Omega^{c:d}&=\ii\epsilon_{cab}\widetilde F^{ab:d},\nonumber\\
 \widetilde G^{ab:de}&=\tfrac12(\widetilde F^{ab:de}+\widetilde F^{ba:de}),
 &\widetilde\Omega^{c:de}&=\ii\epsilon_{cab}\widetilde F^{ab:de}.
 \label{eq:geometry-derivatives}
\end{align}
The metric and curvature derivatives are obtained without suppressing band indices:
\begin{align}
 \widetilde G_{mn}^{ab:d}&=\tfrac12(\widetilde F_{mn}^{ab:d}+\widetilde F_{mn}^{ba:d}),
 &\widetilde\Omega_{mn}^{c:d}&=\ii\epsilon_{cab}\widetilde F_{mn}^{ab:d},\nonumber\\
 \widetilde G_{mn}^{ab:de}&=\tfrac12(\widetilde F_{mn}^{ab:de}+\widetilde F_{mn}^{ba:de}),
 &\widetilde\Omega_{mn}^{c:de}&=\ii\epsilon_{cab}\widetilde F_{mn}^{ab:de}.
 \label{eq:geometry-derivative-components}
\end{align}
For example, the metric dipole of the active subspace is
$\partial_d\sum_{n\in A}\widetilde G_{nn}^{ab}
 =\sum_{n\in A}\widetilde G_{nn}^{ab:d}$.

Taking $\Tr_A$ gives the ordinary first or second derivative of the
corresponding subspace trace.
The same product rule differentiates Eqs.~\eqref{eq:Gweighted} and
\eqref{eq:Hinterpolation}; derivatives of every Hamiltonian factor are required.

\section{BCP and energy-denominator derivatives}
\label{sec:bcp}

\subsection{A symmetric block construction}
For off-block elements ($n\in A$, $l\in B$), define $M$ by Eq.~\eqref{eq:M} and
\begin{equation}
 N_{nl}^a=\frac{M_{nl}^a}{\Delta_{nl}},\qquad
 (M_{nl}^a)^*=M_{ln}^a,\qquad (N_{nl}^a)^*=-N_{ln}^a.
 \label{eq:N}
\end{equation}
The finite-model energy-weighted tensor and BCP are
\begin{align}
 \widetilde{\mathcal F}^{ab}
 &=\tfrac12(N_{AB}^a M_{BA}^b-M_{AB}^a N_{BA}^b),\label{eq:weightedF}\\
 \widetilde{\mathcal{G}}^{ab}
 &=\tfrac12(\widetilde{\mathcal F}^{ab}+\widetilde{\mathcal F}^{ba}).
 \label{eq:bcp-block}
\end{align}
In matrix elements,
\begin{align}
 \widetilde{\mathcal F}_{mn}^{ab}
 &=\tfrac12(N_{ml'}^aM_{l'n}^b-M_{ml'}^aN_{l'n}^b),\nonumber\\
 \widetilde{\mathcal{G}}_{mn}^{ab}
 &=\tfrac12(\widetilde{\mathcal F}_{mn}^{ab}+\widetilde{\mathcal F}_{mn}^{ba}).
 \label{eq:bcp-product-components}
\end{align}

In the energy eigenbasis,
\begin{equation}
 \widetilde{\mathcal F}_{mn}^{ab}
 =\frac12\sum_{l\in B}M_{ml}^a M_{ln}^b
 \left(\frac1{\Delta_{ml}}+\frac1{\Delta_{nl}}\right).
 \label{eq:weightedF-components}
\end{equation}
For a single active band and a complete model complement,
$\widetilde{\mathcal{G}}_{nn}^{ab}=\mathcal{G}_n^{ab}$.
A calculation using only $M$ within a truncated Wannier space does not restore
energy-denominator contributions from all excluded physical bands.

Let $T^{ab}=N_{AB}^aM_{BA}^b-M_{AB}^aN_{BA}^b$.
Hermiticity gives $(T^{ab})^\dagger=T^{ba}$, so the code-friendly form is
\begin{equation}
 \widetilde{\mathcal{G}}^{ab}=\tfrac14[T^{ab}+(T^{ab})^\dagger].
 \label{eq:bcp-code}
\end{equation}
Here the Hermitian symmetrization is explicitly
\begin{equation}
 T_{mn}^{ab}=N_{ml'}^aM_{l'n}^b-M_{ml'}^aN_{l'n}^b,\qquad
 \widetilde{\mathcal{G}}_{mn}^{ab}
 =\tfrac14[T_{mn}^{ab}+(T_{nm}^{ab})^*].
 \label{eq:bcp-code-components}
\end{equation}

This explains the division by four after band-index Hermitian symmetrization
in the archived \code{BCP} class.

\subsection{Connection derivatives and the general block rule}
The first derivative of the off-block coefficient connection is
\begin{equation}
    D_{ln}^{a:d}=\frac{1}{\varepsilon_n-\varepsilon_l}
 \big[\mathcal H_{ln}^{,ad} +\mathcal H_{lp'}^{,a}D_{p'n}^d-D_{lm'}^d\mathcal H_{m'n}^{,a} + \mathcal H_{lp'}^{,d}D_{p'n}^a-D_{lm'}^a\mathcal H_{m'n}^{,d}\big], \label{eq:Dfirst}
\end{equation}
where $m' \in A$ and $p' \in B$ are summed.
Also $M^{a:d}=\bar A^{a:d}+\ii D^{a:d}$, with the corresponding relation at second order.

To differentiate an energy denominator covariantly, define the Sylvester operator
\begin{equation}
 \mathscr L_{\alpha\beta}(Y)=\mathcal H_{\alpha\alpha}Y-Y\mathcal H_{\beta\beta}.
 \label{eq:sylvester}
\end{equation}
Here $\alpha\ne\beta$, and $Y,R,N^a,M^a$ in the Sylvester equations denote
$\alpha\beta$ blocks. Its inverse exists for separated spectra. Equation~\eqref{eq:N} is the
Hamiltonian-gauge evaluation of $\mathscr L_{\alpha\beta}(N^a)=M^a$.
For $\mathscr L_{\alpha\beta}(Y)=R$, product differentiation gives
\begin{align}
 \mathscr L_{\alpha\beta}(Y^{:d})
 &=R^{:d}-\mathcal H_{\alpha\alpha}^{:d}Y+Y\mathcal H_{\beta\beta}^{:d},\label{eq:sylvester1}\\
 \mathscr L_{\alpha\beta}(Y^{:de})
 &=R^{:de}-\mathcal H_{\alpha\alpha}^{:de}Y
 -\mathcal H_{\alpha\alpha}^{:d}Y^{:e}-\mathcal H_{\alpha\alpha}^{:e}Y^{:d}
 +Y^{:d}\mathcal H_{\beta\beta}^{:e}+Y^{:e}\mathcal H_{\beta\beta}^{:d}
       +Y\mathcal H_{\beta\beta}^{:de}.
 \label{eq:sylvester2}
\end{align}
For the $BA$ elements, the Sylvester equation and both derivatives read
\begin{align}
 \mathcal H_{lp'}Y_{p'n}-Y_{lm'}\mathcal H_{m'n}&=R_{ln},\nonumber\\
 \mathcal H_{lp'}Y_{p'n}^{:d}-Y_{lm'}^{:d}\mathcal H_{m'n}
 &=R_{ln}^{:d}-\mathcal H_{lp'}^{:d}Y_{p'n}+Y_{lm'}\mathcal H_{m'n}^{:d},\nonumber\\
 \mathcal H_{lp'}Y_{p'n}^{:de}-Y_{lm'}^{:de}\mathcal H_{m'n}
 &=R_{ln}^{:de}-\mathcal H_{lp'}^{:de}Y_{p'n}
   -\mathcal H_{lp'}^{:d}Y_{p'n}^{:e}
   -\mathcal H_{lp'}^{:e}Y_{p'n}^{:d}\nonumber\\
 &\quad+Y_{lm'}^{:d}\mathcal H_{m'n}^{:e}
   +Y_{lm'}^{:e}\mathcal H_{m'n}^{:d}
   +Y_{lm'}\mathcal H_{m'n}^{:de}.
 \label{eq:sylvester-components}
\end{align}
Thus the sums over $p'\in B$ and $m'\in A$ stay distinct even after
repeated differentiation. The $AB$ equation has elements
\begin{equation}
 \mathcal H_{mn'}Y_{n'l}-Y_{ml'}\mathcal H_{l'l}=R_{ml};
 \label{eq:sylvester-ab-components}
\end{equation}
its derivatives follow the same product rule, with these index spaces.
For BCP use $Y_{ln}=N_{ln}^a$, $R_{ln}=M_{ln}^a$;
for the coefficient connection use $Y_{ln}=D_{ln}^a$,
$R_{ln}=-\mathcal H_{ln}^{,a}$.
In the Hamiltonian gauge the left side of the first line becomes
$(\varepsilon_l-\varepsilon_n)Y_{ln}$.

Set $(Y,R)=(N^a,M^a)$ for BCP denominators, or
$(Y,R)=(D^a_{BA},-\mathcal H_{BA}^{,a})$ for coefficient connections.
These identities preserve intrablock matrix products and supply a general
reference for auditing the multiband implementation.

\subsection{Isolated-band quotient formulae}
For individual nondegenerate eigenlines, use the coefficient-space
phase-covariant derivative $;d$ defined in
Eq.~\eqref{eq:abelian-coefficient-derivative}.
It differs from a multiband block derivative and, in a moving Wannier basis,
from the physical Bloch derivative in Eq.~\eqref{eq:abelian-physical-derivative}.
Let $\Delta_{jk}^{,d}=E_j^d-E_k^d$ and
$\Delta_{jk}^{,de}=\partial_e\partial_d(\varepsilon_j-\varepsilon_k)$.
The scalar comma notation here is simply
$\Delta_{jk}^{,d}=\partial_d\Delta_{jk}$ (and similarly at second order),
with $j\ne k$ and nonzero $\Delta_{jk}$.
Then
\begin{align}
 N_{jk}^{a;d}
 &=\frac{M_{jk}^{a;d}}{\Delta_{jk}}
   -\frac{M_{jk}^{a}\Delta_{jk}^{,d}}{\Delta_{jk}^2},\label{eq:Nfirst}\\
 N_{jk}^{a;de}
 &=\frac{M_{jk}^{a;de}}{\Delta_{jk}}
   -\frac{M_{jk}^{a;d}\Delta_{jk}^{,e}+M_{jk}^{a;e}\Delta_{jk}^{,d}
           +M_{jk}^a\Delta_{jk}^{,de}}{\Delta_{jk}^2}
   +\frac{2M_{jk}^a\Delta_{jk}^{,d}\Delta_{jk}^{,e}}{\Delta_{jk}^3}.
 \label{eq:Nsecond}
\end{align}
These are the compact quotient formulae in \code{Wannier.tex}.
They must not replace Eqs.~\eqref{eq:sylvester1}--\eqref{eq:sylvester2}
when the numerator is evaluated with a multiband colon derivative.
The archived \code{DerBCP} mixes these two prescriptions: it uses the
multiband derivative of $M$, but only diagonal energy gradients in the
denominator derivative. The missing terms are given explicitly below.

\subsection{Missing denominator terms in the archived DerBCP}
\label{sec:bcp-missing-denominator}
\paragraph{The exact matrix-element rule.}
Evaluate Eq.~\eqref{eq:sylvester1} in the Hamiltonian gauge, after taking
the covariant derivative. For $m,n\in A$ and $l,p\in B$, it gives
\begin{align}
 N_{ml}^{a:d}
 &=\frac{M_{ml}^{a:d}
   -\mathcal H_{mn'}^{:d}N_{n'l}^a
   +N_{ml'}^a\mathcal H_{l'l}^{:d}}{\Delta_{ml}},\nonumber\\
 N_{ln}^{a:d}
 &=\frac{M_{ln}^{a:d}
   -\mathcal H_{lp'}^{:d}N_{p'n}^a
   +N_{lm'}^a\mathcal H_{m'n}^{:d}}{\Delta_{ln}}.
 \label{eq:Nfirst-full-components}
\end{align}
Primed indices are implicitly summed in their assigned spaces. Thus the
first line contains both an $A$-space matrix product and a $B$-space matrix
product, even when the unperturbed Hamiltonian is diagonal.

By contrast, extracting only the diagonal entries of the Hamiltonian
derivative, as in the archived \code{Vn}, \code{Vl}, and \code{deltaV}
operations, gives the estimates
\begin{align}
 \widehat N_{ml}^{a:d}
 &=\frac{M_{ml}^{a:d}}{\Delta_{ml}}
   -\frac{N_{ml}^a(E_m^d-E_l^d)}{\Delta_{ml}},\nonumber\\
 \widehat N_{ln}^{a:d}
 &=\frac{M_{ln}^{a:d}}{\Delta_{ln}}
   -\frac{N_{ln}^a(E_l^d-E_n^d)}{\Delta_{ln}},
 \qquad E_j^d=\partial_d\varepsilon_j.
 \label{eq:Nfirst-diagonal-estimate}
\end{align}
The hat denotes an incomplete estimate of the colon derivative, not a
different gauge choice. The scalar quotient rule
\eqref{eq:Nfirst} remains correct for consistently evaluated Abelian
semicolon derivatives; the error is applying its denominator term to a
multiband colon derivative of the numerator.

To expose the omission without suppressing any implicit sums, define the
off-diagonal parts within each space, in the Hamiltonian gauge,
\begin{equation}
 \mathcal H_{mn,\mathrm{off}}^{:d}
 =\mathcal H_{mn}^{:d}-\delta_{mn}E_m^d,\qquad
 \mathcal H_{lp,\mathrm{off}}^{:d}
 =\mathcal H_{lp}^{:d}-\delta_{lp}E_l^d.
 \label{eq:intraspace-Hder-off}
\end{equation}
There is no sum over the unprimed indices in this definition. The omitted
contributions $\delta N^{a:d}=N^{a:d}-\widehat N^{a:d}$ are exactly
\begin{align}
 \delta N_{ml}^{a:d}
 &=\frac{-\mathcal H_{mn',\mathrm{off}}^{:d}N_{n'l}^a
          +N_{ml'}^a\mathcal H_{l'l,\mathrm{off}}^{:d}}{\Delta_{ml}},\nonumber\\
 \delta N_{ln}^{a:d}
 &=\frac{-\mathcal H_{lp',\mathrm{off}}^{:d}N_{p'n}^a
          +N_{lm'}^a\mathcal H_{m'n,\mathrm{off}}^{:d}}{\Delta_{ln}}.
 \label{eq:Nfirst-missing-components}
\end{align}
For a smooth nondegenerate eigenframe, the colon definition itself shows
why these terms need not vanish:
\begin{equation}
 \mathcal H_{mn}^{:d}
 =\delta_{mn}E_m^d+(\varepsilon_n-\varepsilon_m)D_{mn}^d
 =\mathcal H_{mn}^{,d}.
 \label{eq:Hdiag-not-Hderdiag}
\end{equation}
Diagonalizing $\mathcal H$ therefore does not diagonalize its covariant
derivative. A parallel-transport frame can remove the intragroup connection
at a point, but then the ordinary derivative of the Hamiltonian matrix
retains the same off-diagonal information.

Keeping the numerator derivatives fixed, the resulting correction to the
BCP derivative is
\begin{align}
 \delta T_{mn}^{ab:d}
 &=\delta N_{ml'}^{a:d}M_{l'n}^b
   -M_{ml'}^a\delta N_{l'n}^{b:d},\nonumber\\
 \delta\widetilde{\mathcal G}_{mn}^{ab:d}
 &=\tfrac14\big[\delta T_{mn}^{ab:d}
                    +(\delta T_{nm}^{ab:d})^*\big].
 \label{eq:DerBCP-missing-components}
\end{align}
Hermitian symmetrization does not in general remove this correction, nor
does taking the $A$-space trace. If both $A$ and $B$ contain just one band,
the omitted sums vanish; this explains why a two-band Haldane check alone
does not detect the error. Choosing a single active band $A=\{n\}$ is
insufficient in a larger model: the off-diagonal $B$-space terms remain.
More generally, the diagonal estimate is sufficient only when the
correction in Eq.~\eqref{eq:Nfirst-missing-components} vanishes for the
chosen matrix elements and derivative direction.

\paragraph{Why the Sylvester operator resolves the omission.}
The energy denominator is the eigenbasis representation of an inverse
linear operator on matrices, $N^a=\mathscr L_{\alpha\beta}^{-1}(M^a)$.
Its covariant derivative acts on both the input and the operator:
\begin{align}
 N^{a:d}
 &=\mathscr L_{\alpha\beta}^{-1}
    \big[M^{a:d}-\mathscr L_{\alpha\beta}^{:d}(N^a)\big],\nonumber\\
 \mathscr L_{\alpha\beta}^{:d}(Y)
 &\equiv\mathcal H_{\alpha\alpha}^{:d}Y
            -Y\mathcal H_{\beta\beta}^{:d}.
 \label{eq:sylvester-inverse-derivative}
\end{align}
Here $\mathscr L^{:d}$ differentiates the operator coefficients while
holding its argument fixed. Equivalently,
$(\mathscr L^{-1})^{:d}=-\mathscr L^{-1}\mathscr L^{:d}\mathscr L^{-1}$.
The two Hamiltonian-derivative products produce all the terms in
Eq.~\eqref{eq:Nfirst-full-components}; replacing them by diagonal energy
gradients discards precisely Eq.~\eqref{eq:Nfirst-missing-components}.

This formulation is also covariant under independent rotations within
$A$ and $B$. With $\mathcal H'_{\alpha\alpha}
=W_\alpha^\dagger\mathcal H_{\alpha\alpha}W_\alpha$,
\begin{equation}
 \mathscr L'_{\alpha\beta}(W_\alpha^\dagger YW_\beta)
 =W_\alpha^\dagger\mathscr L_{\alpha\beta}(Y)W_\beta.
 \label{eq:sylvester-covariance}
\end{equation}
For example, $Y'_{ln}=(W^\dagger)_{ll'}Y_{l'm'}W_{m'n}$ and the right-hand
side $R_{ln}$ transforms by the same matrix-element rule. Uniqueness of
the solution then preserves this covariance for $N$ and its colon
derivatives. Inverting $\mathscr L$ requires disjoint $A$ and $B$ spectra,
not nondegeneracy within either space; a vanishing cross-space gap must
instead be treated by regrouping the bands or specifying another model.

In practice no large Sylvester solver is necessary in the Hamiltonian
gauge: form the complete right-hand side with ordinary matrix products,
then divide its $ml$ or $ln$ element by the corresponding cross-space
energy difference. This is the rule used by
\code{calculators_collect/formula/covariant_basic.py}.
Differentiating the same operator equation again yields
Eq.~\eqref{eq:sylvester2}, including both mixed first-derivative products
and the second Hamiltonian derivative. Thus the same construction fixes
the denominator treatment in \code{DerBCP} and supplies a consistent
one for \code{Der2BCP}.

\subsection{BCP dipole and quadrupole kernels}
With either consistently evaluated block derivatives or isolated-band derivatives,
\begin{align}
 \widetilde{\mathcal F}^{ab:d}
 &=\tfrac12\big[
 N_{AB}^{a:d}M_{BA}^b+N_{AB}^aM_{BA}^{b:d}
 -M_{AB}^{a:d}N_{BA}^b-M_{AB}^aN_{BA}^{b:d}\big],\label{eq:bcp-first}\\
 \widetilde{\mathcal F}^{ab:de}
 &=\tfrac12\big[\mathscr P^{de}(N_{AB}^a,M_{BA}^b)
                  -\mathscr P^{de}(M_{AB}^a,N_{BA}^b)\big],\label{eq:bcp-second}\\
 \widetilde{\mathcal{G}}^{ab:d}
 &=\tfrac12(\widetilde{\mathcal F}^{ab:d}
             +\widetilde{\mathcal F}^{ba:d}),\nonumber\\
 \widetilde{\mathcal{G}}^{ab:de}
 &=\tfrac12(\widetilde{\mathcal F}^{ab:de}
             +\widetilde{\mathcal F}^{ba:de}).
 \label{eq:bcp-sym}
\end{align}
The first and second derivatives in matrix-element form are
\begin{align}
 \widetilde{\mathcal F}_{mn}^{ab:d}
 &=\tfrac12\big[
 N_{ml'}^{a:d}M_{l'n}^b+N_{ml'}^aM_{l'n}^{b:d}
 -M_{ml'}^{a:d}N_{l'n}^b-M_{ml'}^aN_{l'n}^{b:d}\big],\nonumber\\
 \widetilde{\mathcal F}_{mn}^{ab:de}
 &=\tfrac12\big[
 N_{ml'}^{a:de}M_{l'n}^b+N_{ml'}^{a:d}M_{l'n}^{b:e}
 +N_{ml'}^{a:e}M_{l'n}^{b:d}+N_{ml'}^aM_{l'n}^{b:de}\nonumber\\
 &\qquad-M_{ml'}^{a:de}N_{l'n}^b-M_{ml'}^{a:d}N_{l'n}^{b:e}
 -M_{ml'}^{a:e}N_{l'n}^{b:d}-M_{ml'}^aN_{l'n}^{b:de}\big],\nonumber\\
 \widetilde{\mathcal{G}}_{mn}^{ab:d}
 &=\tfrac12(\widetilde{\mathcal F}_{mn}^{ab:d}+\widetilde{\mathcal F}_{mn}^{ba:d}),\nonumber\\
 \widetilde{\mathcal{G}}_{mn}^{ab:de}
 &=\tfrac12(\widetilde{\mathcal F}_{mn}^{ab:de}+\widetilde{\mathcal F}_{mn}^{ba:de}).
 \label{eq:bcp-derivative-components}
\end{align}

The two factors of $1/2$ in this construction are independent.
For a smooth subspace, $\Tr_A\widetilde{\mathcal{G}}^{ab:d}$ and
$\Tr_A\widetilde{\mathcal{G}}^{ab:de}$ are the dipole and quadrupole
geometric kernels before occupations and response prefactors are applied.

\section{Explicit band-sum and matrix formulae from QGD}
\label{sec:qgd}

\subsection{Domain and matrix ingredients}
This section applies to a smooth Hamiltonian in a fixed orthonormal basis,
with nondegenerate bands and nonzero denominators.
It gives an independent representation of the internal geometry.
Define energy-derivative matrices, rather than physical velocities,
\begin{align}
 V_{jk}^a&=\langle j|\partial_a\mathcal H|k\rangle,
 &\Gamma_{jk}^{ab}&=\langle j|\partial_a\partial_b\mathcal H|k\rangle,
 \nonumber\\
 (K_q^a)_{jk}&=\begin{cases}V_{jk}^a/\Delta_{jk}^q,&j\ne k,\\0,&j=k,\end{cases}
 &(V_F^a)_{jk}&=\delta_{jk}E_j^a.
 \label{eq:qgd-ingredients}
\end{align}
Division by $\Delta^q$ is elementwise and excludes the diagonal.
Here $q$ is a positive integer power, not a band or Cartesian index.
Products of $K_q$ are ordinary matrix products over Full.
In this convention $D_{jk}^a=-(K_1^a)_{jk}$ and
$A_{jk}^a=-\ii(K_1^a)_{jk}$ for $j\ne k$.
Thus $K_1$ alone is not the interband Berry connection.

The QGT can be evaluated by either expression
\begin{equation}
 Q_n^{ab}=\sum_{j\ne n}\frac{V_{nj}^aV_{jn}^b}{\Delta_{nj}^2}
          =-[K_1^aK_1^b]_{nn}.
 \label{eq:qgt-matrix}
\end{equation}
Only the diagonal of the matrix product is asserted to equal the
single-band QGT. Its off-diagonal entries are not the general projected
matrix of Eq.~\eqref{eq:Fprojected}.

\subsection{QGD-I and QGD-II}
Write $\partial_cQ_n^{ab}=\mathcal I_n^{ab;c}+\mathcal J_n^{ab;c}$.
The one-sum contribution is
\begin{equation}
 \mathcal I_n^{ab;c}=\sum_{j\ne n}\left[
 \frac{\Gamma_{nj}^{ca}V_{jn}^b+V_{nj}^a\Gamma_{jn}^{cb}}{\Delta_{nj}^2}
 -\frac{2V_{nj}^aV_{jn}^b(E_n^c-E_j^c)}{\Delta_{nj}^3}\right].
 \label{eq:qgd-I}
\end{equation}
The two-sum contribution is
\begin{align}
 \mathcal J_n^{ab;c}=\sum_{j\ne n}\bigg[
 &\sum_{k\ne n}\frac{V_{nk}^cV_{kj}^aV_{jn}^b
                        +V_{nj}^aV_{jk}^bV_{kn}^c}
                       {\Delta_{nj}^2\Delta_{nk}}
 \nonumber\\
 +&\sum_{k\ne j}\frac{V_{nk}^aV_{kj}^cV_{jn}^b
                        +V_{nj}^aV_{jk}^cV_{kn}^b}
                       {\Delta_{nj}^2\Delta_{jk}}\bigg].
 \label{eq:qgd-II}
\end{align}
The corresponding compact matrix expressions are
\begin{align}
 \mathcal I_n^{ab;c}
 &=\big[\Gamma^{ca}K_2^b+K_2^a\Gamma^{cb}
        -2[V_F^c,K_1^a]K_2^b\big]_{nn},\label{eq:qgd-I-matrix}\\
 \mathcal J_n^{ab;c}
 &=\big[[K_1^c,V^a]K_2^b+K_2^a[K_1^c,V^b]\big]_{nn}.
 \label{eq:qgd-II-matrix}
\end{align}
No real part is taken until selecting the desired observable:
\begin{equation}
 \partial_cG_n^{ab}=\Rea\partial_cQ_n^{ab},\qquad
 \partial_c\Omega_n^{ab}=-2\Ima\partial_cQ_n^{ab}.
 \label{eq:qgd-observables}
\end{equation}
Equations~\eqref{eq:qgd-I}--\eqref{eq:qgd-II-matrix} retain the
one-sum/two-sum distinction in \code{QGD.tex} without repeating the intermediate bra-ket expansions.

\subsection{A complete quotient identity for the second derivative}
The QGQ expansions in the source notes are incomplete.
For a compact check of any completed expansion, define the phase-invariant pair
$Z_{nj}^{ab}=V_{nj}^aV_{jn}^b$ for $j\ne n$.
Differentiating $Q_n^{ab}=\sum_{j\ne n}Z_{nj}^{ab}/\Delta_{nj}^2$ twice gives
\begin{align}
 \partial_d\partial_c Q_n^{ab}=\sum_{j\ne n}\bigg[
 &\frac{\partial_d\partial_c Z_{nj}^{ab}}{\Delta_{nj}^2}
 -\frac{2\big[(\partial_c Z_{nj}^{ab})\Delta_{nj}^{,d}
              +(\partial_d Z_{nj}^{ab})\Delta_{nj}^{,c}
              +Z_{nj}^{ab}\Delta_{nj}^{,cd}\big]}{\Delta_{nj}^3}
 \nonumber\\
 &+\frac{6Z_{nj}^{ab}\Delta_{nj}^{,c}\Delta_{nj}^{,d}}{\Delta_{nj}^4}\bigg].
 \label{eq:qgq-quotient}
\end{align}
In particular, the last coefficient is $+6$ in this scalar band-sum convention.
The numerator derivatives include changes of both eigenstates and operators.
For an arbitrary operator matrix $X$, they can be generated using
\begin{equation}
 X_{jk}^{;c}=\langle j|\partial_c\hat X|k\rangle+[K_1^c,X]_{jk},\qquad
 (K_q^a)_{jk}^{;c}=\frac{(V^a)_{jk}^{;c}}{\Delta_{jk}^q}
 -q\frac{V_{jk}^a\Delta_{jk}^{,c}}{\Delta_{jk}^{q+1}},\quad j\ne k.
 \label{eq:qgd-recursion}
\end{equation}
Here $;c$ removes only eigenline phase connections, with the same choice as in
Sec.~\ref{sec:bcp}. Differentiating $\Gamma$ includes
$\langle j|\partial_c\partial_a\partial_b\mathcal H|k\rangle$.
It cannot be omitted for a general tight-binding Hamiltonian.
This identity is a complete differentiation rule, not a claim that the
partial QGQ matrix expansions in the source notes are complete.

\section{Energy derivatives and Drude building blocks}
\label{sec:energy}
Define the rotated bare derivatives
\begin{equation}
 V^a=\mathcal H^{,a},\quad W^{ab}=\mathcal H^{,ab},\quad
 Y^{abc}=\mathcal H^{,abc},\quad Z^{abcd}=\mathcal H^{,abcd}.
 \label{eq:energy-ingredients}
\end{equation}
Here $Z^{abcd}$ is a four-derivative Hamiltonian matrix, distinct from the
band-pair numerator $Z_{nj}^{ab}$ in Sec.~\ref{sec:qgd}.
For example,
\begin{align}
 V_{AA}^{a:b}&=W_{AA}^{ab}+V_{AB}^aD_{BA}^b-D_{AB}^bV_{BA}^a,
 \label{eq:invmass}\\
 W_{AA}^{ab:c}&=Y_{AA}^{abc}+W_{AB}^{ab}D_{BA}^c-D_{AB}^cW_{BA}^{ab},
 \label{eq:derW}\\
 V_{AA}^{a:bc}&=W_{AA}^{ab:c}
 +V_{AB}^{a:c}D_{BA}^b+V_{AB}^aD_{BA}^{b:c}
 -D_{AB}^{b:c}V_{BA}^a-D_{AB}^bV_{BA}^{a:c}.
 \label{eq:third-energy}
\end{align}
A further product derivative yields
\begin{align}
 V_{AA}^{a:bcd}={}&W_{AA}^{ab:cd}
 +V_{AB}^{a:cd}D_{BA}^b+V_{AB}^{a:c}D_{BA}^{b:d}
 +V_{AB}^{a:d}D_{BA}^{b:c}+V_{AB}^aD_{BA}^{b:cd}
 \nonumber\\
 &-D_{AB}^{b:cd}V_{BA}^a-D_{AB}^{b:c}V_{BA}^{a:d}
 -D_{AB}^{b:d}V_{BA}^{a:c}-D_{AB}^bV_{BA}^{a:cd}.
 \label{eq:fourth-energy}
\end{align}
The same derivative chain, with free band indices $m,n\in A$, is
\begin{align}
 V_{mn}^{a:b}&=W_{mn}^{ab}+V_{ml'}^aD_{l'n}^b-D_{ml'}^bV_{l'n}^a,\nonumber\\
 W_{mn}^{ab:c}&=Y_{mn}^{abc}+W_{ml'}^{ab}D_{l'n}^c-D_{ml'}^cW_{l'n}^{ab},\nonumber\\
 V_{mn}^{a:bc}&=W_{mn}^{ab:c}
 +V_{ml'}^{a:c}D_{l'n}^b+V_{ml'}^aD_{l'n}^{b:c}
 -D_{ml'}^{b:c}V_{l'n}^a-D_{ml'}^bV_{l'n}^{a:c},\nonumber\\
 V_{mn}^{a:bcd}&=W_{mn}^{ab:cd}
 +V_{ml'}^{a:cd}D_{l'n}^b+V_{ml'}^{a:c}D_{l'n}^{b:d}
 +V_{ml'}^{a:d}D_{l'n}^{b:c}+V_{ml'}^aD_{l'n}^{b:cd}\nonumber\\
 &\quad-D_{ml'}^{b:cd}V_{l'n}^a-D_{ml'}^{b:c}V_{l'n}^{a:d}
 -D_{ml'}^{b:d}V_{l'n}^{a:c}-D_{ml'}^bV_{l'n}^{a:cd}.
 \label{eq:energy-derivative-components}
\end{align}

For a smooth isolated group,
\begin{equation}
 \partial_a\partial_b\sum_{n\in A}\varepsilon_n=\Tr_A V^{a:b},\quad
 \partial_a\partial_b\partial_c\sum_{n\in A}\varepsilon_n=\Tr_A V^{a:bc},\quad
 \partial_a\partial_b\partial_c\partial_d\sum_{n\in A}\varepsilon_n=\Tr_A V^{a:bcd}.
 \label{eq:energy-traces}
\end{equation}
Each trace is simply the sum of diagonal matrix elements:
\begin{equation}
 \Tr_A V^{a:b}=\sum_{n\in A}V_{nn}^{a:b},\qquad
 \Tr_A V^{a:bc}=\sum_{n\in A}V_{nn}^{a:bc},\qquad
 \Tr_A V^{a:bcd}=\sum_{n\in A}V_{nn}^{a:bcd}.
 \label{eq:energy-trace-components}
\end{equation}

For a single isolated band the inverse effective mass is
$(m_n^{-1})^{ab}=\hbar^{-2}\partial_a\partial_b\varepsilon_n$.
The operator Hessian $W_{nn}^{ab}$ differs from the band-energy Hessian:
\begin{equation}
 \partial_a\partial_b\varepsilon_n
 =W_{nn}^{ab}+2\Rea\sum_{j\ne n}\frac{V_{nj}^aV_{jn}^b}{\Delta_{nj}}.
 \label{eq:hessian-distinction}
\end{equation}
Thus the orbital-quadrupole derivative of an operator Hessian must not be
replaced by \code{Der3E}. They correspond to different objects.

\section{Orbital magnetic quantities}
\label{sec:orbital}

\subsection{Orbital moment and zero-temperature magnetization}
With the charge convention in Eq.~\eqref{eq:conventions},
\begin{align}
 \widetilde m^c&=\frac{\ii e}{2\hbar}\epsilon_{cab}
                      (\widetilde{\mathcal E}^{ab}-\widetilde H^{ab}),\label{eq:moment}\\
 \widetilde{\mathcal M}^c(\mu)
 &=-\frac{\ii e}{2\hbar}\epsilon_{cab}
    (\widetilde H^{ab}+\widetilde{\mathcal E}^{ab}-2\mu\widetilde F^{ab}).
 \label{eq:magnetization}
\end{align}
The corresponding band-resolved matrix elements are
\begin{align}
 \widetilde m_{mn}^c&=\frac{\ii e}{2\hbar}\epsilon_{cab}
 (\widetilde{\mathcal E}_{mn}^{ab}-\widetilde H_{mn}^{ab}),\nonumber\\
 \widetilde{\mathcal M}_{mn}^c(\mu)&=-\frac{\ii e}{2\hbar}\epsilon_{cab}
 (\widetilde H_{mn}^{ab}+\widetilde{\mathcal E}_{mn}^{ab}-2\mu\widetilde F_{mn}^{ab}).
 \label{eq:orbital-components}
\end{align}

The first is the symmetrized orbital-moment matrix; for an exactly degenerate
active group it reduces to the usual moment matrix.
For a single band,
$\mathcal M_n^c=m_n^c+(e/\hbar)(\mu-\varepsilon_n)\Omega_n^c$.
At zero temperature, using the occupied projector at each wave vector,
\begin{equation}
 M_{\mathrm{orb}}^c=\int[\dd\bm k]\Tr_{\mathrm{occ}}
                    \widetilde{\mathcal M}^c(\mu).
 \label{eq:magnetization-integral}
\end{equation}
Taking $A$ to be the occupied subspace gives the explicit sum
\begin{equation}
 M_{\mathrm{orb}}^c=\int[\dd\bm k]\sum_{n\in A}
 \widetilde{\mathcal M}_{nn}^c(\mu).
 \label{eq:magnetization-integral-components}
\end{equation}

At fixed $\mu$, the derivative of Eq.~\eqref{eq:magnetization} is
\begin{equation}
 (\widetilde{\mathcal M}^c)^{:d}
 =-\frac{\ii e}{2\hbar}\epsilon_{cab}
 (\widetilde H^{ab:d}+\widetilde{\mathcal E}^{ab:d}-2\mu\widetilde F^{ab:d}).
 \label{eq:magnetization-derivative}
\end{equation}
At fixed $\mu$, its matrix elements satisfy
\begin{equation}
 \widetilde{\mathcal M}_{mn}^{c:d}
 =-\frac{\ii e}{2\hbar}\epsilon_{cab}
 (\widetilde H_{mn}^{ab:d}+\widetilde{\mathcal E}_{mn}^{ab:d}
                         -2\mu\widetilde F_{mn}^{ab:d}).
 \label{eq:magnetization-derivative-components}
\end{equation}

Every projector derivative in $\widetilde H$ and $\widetilde{\mathcal E}$ is included
through the interpolation and product rules. Keeping only
$\langle\partial_a u_m|Q(\partial_d\hat{\mathcal H})Q|\partial_b u_n\rangle$
is insufficient.

\subsection{Orbital magnetic quadrupole: four-term reference}
\label{sec:mqm}
To avoid the duplicated prefactors in \code{Wannier.tex}, this subsection
uses the $e=\hbar=1$ convention of Gao and Xiao~\cite{GaoXiao2018}.
It supplies the reduced kernel for comparison with the archive, not a new SI conversion.
For an isolated band choose $A=\{n\}$ and $B=\mathrm{Full}\setminus\{n\}$, and define
\begin{align}
 \mathfrak m_{ln}^{b}
 &=\frac12\epsilon_{bcd}\sum_{p\in B}
       (v_{lp}^c+v_n^c\delta_{lp})A_{pn}^d,\label{eq:interband-moment}\\
 \Gamma_n^{ad}&=\langle u_n|\partial_a\partial_d\hat{\mathcal H}|u_n\rangle.
\end{align}
Here $v^c=\partial_c\hat{\mathcal H}$ in this subsection's units,
and $G_n^{ab}$ denotes the quantum metric throughout.
Define four band kernels,
\begin{align}
 C_n^{ab}&=\frac23\Rea\sum_{l\in B}A_{nl}^a\mathfrak m_{ln}^b,
 &I_n^{ab}&=-\frac1{12}\epsilon_{cdb}\partial_c\Gamma_n^{ad},
 \nonumber\\
 L_n^{ab}&=-2\Rea\sum_{l\in B}
                 \frac{A_{nl}^a\mathfrak m_{ln}^b}{\Delta_{nl}},
 &K_n^{ab}&=-\frac13\epsilon_{cdb}\partial_cG_n^{ad}.
 \label{eq:mqm-kernels}
\end{align}
For the full-band sum, evaluate each isolated-band kernel above with
$n$ relabeled as $j$, and with its own complementary space
$B=\mathrm{Full}\setminus\{j\}$. The quadrupole density is
\begin{equation}
 \mathcal Q^{ab}=\int[\dd\bm k]\sum_{j\in\mathrm{Full}}
 \left\{f_j(C_j^{ab}+I_j^{ab})+\Phi_j(L_j^{ab}+K_j^{ab})\right\},
 \label{eq:mqm}
\end{equation}
where
\begin{equation}
 f_j=\frac1{1+e^{(\varepsilon_j-\mu)/(k_BT)}},\qquad
 \Phi_j=-k_BT\ln\left(1+e^{(\mu-\varepsilon_j)/(k_BT)}\right).
 \label{eq:grand-potential}
\end{equation}
At $T=0$, $f_j=\Theta(\mu-\varepsilon_j)$ and
$\Phi_j=(\varepsilon_j-\mu)f_j$.
The source-note labels ME1, W, ME2, QM correspond respectively to the
archived v1.8 labels CL, IMD, ME, QMD.

The physical operator Hessian in Eq.~\eqref{eq:mqm-kernels} equals the fixed-basis
matrix $W$ of Sec.~\ref{sec:energy} in that representation.
With a moving Wannier basis, operator derivatives and derivatives of
$\mathcal H^W$ require additional basis-connection terms.
This distinction must be checked when comparing the archived \code{DerWln}
implementation to the physical quadrupole formula.

The archived v1.8 calculator uses $(E-\mu)f$ for ME and QMD and defaults to
\code{constant_factor=1}. Its weighting agrees with
Eq.~\eqref{eq:grand-potential} at zero temperature, and does not by itself
implement the finite-temperature grand potential.
No additional $e/(2\hbar)$ factor should be inserted into
Eq.~\eqref{eq:mqm} while retaining the convention in Eq.~\eqref{eq:interband-moment}.
Conversion back to physical units must be derived for the full expression.

\section{Bridge to the v1.7 and v1.8 calculator catalogue}
\label{sec:catalogue}

\subsection{Fermi-sea and Fermi-surface kernels}
For a smooth periodic band quantity $X_j$ and vanishing boundary terms,
with all band sums over $j\in\mathrm{Full}$,
\begin{align}
 \int[\dd\bm k]\sum_j f_j\partial_aX_j
 &=-\int[\dd\bm k]\sum_j f'_j E_j^a X_j,\label{eq:sea-surface}\\
 \int[\dd\bm k]\sum_j f_j\partial_a\partial_dX_j
 &=-\int[\dd\bm k]\sum_j f'_j E_j^a\partial_dX_j,
 \label{eq:sea-surface2}
\end{align}
where $f'_j=\partial f/\partial\varepsilon|_{\varepsilon_j}$ is negative.
A surface kernel defined with $-f'$ has the same sign as the sea integral.
The code's occupation-derivative convention and outer factor must therefore
be checked together. At band crossings, use smooth group traces and the
matrix occupation rule rather than applying isolated-band formulae blindly.

\subsection{Tensor axes and formula-class index}
The following table records geometric kernel axes, excluding any leading
wave-vector and chemical-potential axes. It is a navigation aid for the
subsequent full version reconciliation.
\begin{center}
\setlength{\tabcolsep}{6pt}
\begin{tabular}{lll}
\toprule
Formula class & Geometric kernel & Raw Cartesian axes\\
\midrule
\code{Metric} & $G^{ab}$ or subspace metric & $(a,b)$\\
\code{DerMetric} & $G^{ab:d}$ & $(a,b,d)$\\
\code{BCP} & $\mathcal{G}^{ab}$ & $(a,b)$\\
\code{DerBCP} & $\mathcal{G}^{ab:d}$ & $(a,b,d)$\\
\code{Der2BCP} & $\mathcal{G}^{ab:de}$ & $(a,b,d,e)$\\
\code{DerOmega}, \code{Der2Omega} & $\Omega^{c:d}$, $\Omega^{c:de}$ & $(c,d)$, $(c,d,e)$\\
\code{VelBCP} & $E^c\mathcal{G}^{ab}$ & $(c,a,b)$\\
\code{VelDerBCP} & $E^c\mathcal{G}^{ab:d}$ & $(c,a,b,d)$\\
\code{VelDerOmega} & $E^a\Omega^{c:d}$ & $(a,c,d)$\\
\code{InvMass}, \code{Der3E}, \code{Der4E} & band-energy derivatives & derivative order\\
\code{Qorb_GaoXiao_*} & Eq.~\eqref{eq:mqm} & $(a,b)$\\
\bottomrule
\end{tabular}
\end{center}
For a common dipole convention $(d,a,b)$, transpose the three axes of
\code{DerBCP} from $(a,b,d)$ to $(d,a,b)$.
For a common quadrupole convention $(d,a,b,e)$, the corresponding permutation
of \code{Der2BCP} is $(2,0,1,3)$ in zero-based tensor-axis notation.
Only the trailing tensor axes should be permuted when leading axes are present.

\subsection{Items for the subsequent complete version reconciliation}
The archived source and \code{CHANGE_REPORT.md} identify the following scope:
\begin{itemize}
 \item v1.7 adds metric, BCP dipole/quadrupole, Berry-curvature quadrupole,
 fourth-energy-derivative, and early orbital-quadrupole calculators.
 \item v1.8 retains several formula classes but does not expose the complete
 v1.7 calculator set. It adds the four-term Gao--Xiao quadrupole and
 \code{Beta_GaoXiao} thermoelectric response.
 \item \code{Der2BCP_test} in v1.7 raises \code{NotImplementedError}.
 The production \code{Der2BCP} also carries a bug-related TODO in the source;
 a class name or a completed symbolic identity is not evidence of numerical validation.
 \item BCP denominator derivatives, active/complement choices, tensor axes,
 occupation signs, and SI prefactors need class-by-class comparison.
 Group-metric traces and sums of isolated-band metrics must be distinguished.
\end{itemize}
The thermoelectric response and the full nonlinear conductivity prefactors
are outside the completed source-note consolidation here.
They should be added as separate catalogue entries with their field convention,
thermal kernel, component symmetrization, and version-specific code mapping.

\appendix
\section{Corrections and retained evidence boundaries}
\label{app:corrections}
\begin{enumerate}
 \item \textbf{Connection and derivative notation.}
 The source notes used the same state notation for Bloch functions and
 columns of $U$. Equations~\eqref{eq:D}--\eqref{eq:M} distinguish them.
 The ordinary derivative $\partial_d X$ was also confused with the
 operator/comma derivative $X^{,d}$; Sec.~\ref{sec:derivative-conventions} repairs that distinction.
 The right-block gauge transformation has left index $n$, not the unrelated index $m$.
 \item \textbf{Curvature and Hermiticity.}
 The anti-Hermitian relation is $(D_{jk}^a)^*=-D_{kj}^a$.
 The $z$-curvature contraction uses $xy-yx$, not $xy-yz$.
 The internal trace curvature has the positive sign in
 Eq.~\eqref{eq:curvature-check}.
 The $+2\Ima Q$ convention in \code{QGD.tex} is replaced consistently by
 Eq.~\eqref{eq:metric-curvature}.
 \item \textbf{BCP factors and denominator derivatives.}
 The metric symmetrization and the symmetric energy-denominator construction
 each contribute a factor of $1/2$.
 An expanded derivative in \code{Wannier.tex} lost one of these factors
 and repeated terms instead of exchanging all $a,b$ indices.
 Equations~\eqref{eq:bcp-first}--\eqref{eq:bcp-sym} retain the two operations explicitly.
 The Sylvester identities supplement the single-band quotient rule without
 assuming diagonal intrablock velocity matrices.
 The archived \code{DerBCP} omits the terms in
 Eq.~\eqref{eq:Nfirst-missing-components} by combining a multiband
 numerator derivative with diagonal energy-gradient differences.
 Section~\ref{sec:bcp-missing-denominator} gives the resulting BCP
 correction and explains the inverse-Sylvester derivative used in the
 corrected \code{calculators_collect} implementation.
 \item \textbf{Incomplete QGQ expansion.}
 \code{Wannier.tex} leaves its part III empty, and \code{QGD.tex} repeats
 the part-I label on a different group of terms.
 These partial expressions are not reproduced as a complete result.
 Equation~\eqref{eq:qgq-quotient} records the full quotient identity, including
 the $+6$ denominator term and the band-energy Hessian difference.
 \item \textbf{Orbital terms.}
 The energy-weighted Wannier matrix uses derivatives $a,b$, not $b,b$.
 Orbital-moment, magnetization, and interband magnetic-kernel symbols are separated.
 The original quadrupole expression mixed an outer $e/(2\hbar)$ with a magnetic
 matrix already containing that factor. Section~\ref{sec:mqm} uses a single
 explicit reduced-unit convention.
 \item \textbf{Exploratory derivations.}
 The unfinished ME1 projector manipulations, partial magnetization derivatives,
 and terminal Green-function identities are retained in the original notes.
 They are not needed for the reusable formula chain in this handbook.
 In particular, differentiation of a projected product must include all
 projector derivatives; differentiating only its explicit Hamiltonian is incomplete.
\end{enumerate}

\section{Consistency checks for later calculator validation}
\label{app:checks}
Useful algebraic checks include
\begin{align}
 Q_n^{ab}&=(Q_n^{ba})^*,\qquad G_n^{aa}\ge0,\nonumber\\
 \widetilde F_{mn}^{ab:d}&=(\widetilde F_{nm}^{ba:d})^*,\qquad
 \widetilde F_{mn}^{ab:de}=(\widetilde F_{nm}^{ba:de})^*,\nonumber\\
 \mathcal{G}_n^{ab}&=\mathcal{G}_n^{ba},\nonumber\\
 \partial_d\partial_cX_n&=\partial_c\partial_dX_n
 \quad\text{for a smooth gauge-invariant scalar }X_n.
\end{align}
Compare the QGD band sums, matrix products, and centered finite differences
on a gapped multiband model. Test the BCP block denominator identities under
independent unitary rotations in both blocks.
For Brillouin-zone integrals, compare sea/surface kernels only after matching
axes, signs, units, occupations, and integration normalization.
These checks do not replace convergence tests on the eventual material model.

% Embedded bibliography keeps this working copy self-contained for the native
% single-file LaTeX compiler. Original notes/ref.bib is deliberately unchanged.
\begin{thebibliography}{9}
\bibitem{Liu2023CovariantderivativesBerry}
X. Liu, S. S. Tsirkin, and I. Souza,
Covariant derivatives of Berry-type quantities: Application to nonlinear transport,
\href{https://arxiv.org/abs/2303.10129}{arXiv:2303.10129} (2023).
\bibitem{GaoPRL2014FieldInducedPositional}
Y. Gao, S. A. Yang, and Q. Niu,
Field Induced Positional Shift of Bloch Electrons and Its Dynamical Implications,
\href{https://doi.org/10.1103/PhysRevLett.112.166601}{Phys. Rev. Lett. \textbf{112}, 166601} (2014).
\bibitem{GaoXiao2018}
Y. Gao and D. Xiao,
Orbital magnetic quadrupole moment and nonlinear anomalous thermoelectric transport,
\href{https://doi.org/10.1103/PhysRevB.98.060402}{Phys. Rev. B \textbf{98}, 060402(R)} (2018).
\end{thebibliography}
\end{document}
